\documentclass[11pt]{article}
\usepackage[margin=1in]{geometry}
\usepackage{times}
\usepackage{amsmath}
\usepackage{enumitem}
\usepackage{parskip}

\setlength{\parskip}{6pt}
\setlength{\parindent}{0pt}

\title{\vspace{-1.5cm}\textbf{Research Proposal: Context-Aware Interpretation\\Thresholds for PABAK and Kappa-Type Metrics}}
\author{}
\date{}

\begin{document}
\maketitle
\vspace{-1.2cm}

\section*{The Problem}

Researchers routinely label agreement coefficients ($\kappa$, PABAK) using qualitative scales (e.g., Landis \& Koch, 1977: ``slight,'' ``moderate,'' ``substantial''). These labels have two fundamental issues:

\begin{itemize}[nosep, leftmargin=1.5em]
    \item They were never empirically calibrated. They are heuristics from the 1970s applied as if they were universal standards.
    \item They are borrowed across metrics that measure different things. PABAK uses a uniform chance model (Brennan \& Prediger, 1981; Byrt et al., 1993); Cohen's $\kappa$ (Cohen, 1960) uses marginal-based chance. Applying the same labels to both ignores this structural difference.
\end{itemize}

For binary outcomes, PABAK $= 2P_0 - 1$. It is just a linear rescaling of percent agreement. Any fixed ``high/medium/low'' ladder for PABAK is therefore an implicit ladder for percent agreement. The same PABAK value can correspond to very different error profiles depending on prevalence and rater bias (Feinstein \& Cicchetti, 1990; Hoehler, 2000), yet current practice assigns a single label regardless of context (Chen et al., 2009).

\section*{What We Propose}

Replace static interpretation cut points with \textbf{context-conditioned threshold functions} whose meaning is stable across study designs and decision settings.

\subsection*{Research Questions}
\begin{enumerate}[nosep, leftmargin=1.5em]
    \item When does PABAK accurately reflect decision-relevant agreement, and when does it obscure meaningful disagreement?
    \item How often do fixed adjective scales produce misleading labels relative to decision-analytic criteria?
    \item Can threshold functions conditioned on prevalence, bias, $K$, and cost structure outperform static ladders?
\end{enumerate}

\subsection*{Method: Monte Carlo Simulation}
\begin{itemize}[nosep, leftmargin=1.5em]
    \item Generate binary and multi-class rating data from known rater-accuracy models (sensitivity/specificity, misclassification matrices; Hoehler, 2000; Vanacore \& Pellegrino, 2022). Ground truth is defined by operating characteristics and decision loss.
    \item Vary conditions systematically across a full factorial grid:
    \begin{itemize}[nosep, leftmargin=1em]
        \item Prevalence: $p \in \{0.01, 0.05, 0.10, 0.20, 0.50, 0.80, 0.90, 0.95, 0.99\}$
        \item Rater bias: multiple levels of marginal asymmetry
        \item Categories: $K \in \{2, 3, 5, 7\}$, nominal and ordinal
        \item Raters: $R \in \{2, 3, 5, 10\}$
        \item Sample size: $N \in \{30, 50, 100, 200, 500, 1000\}$
    \end{itemize}
    \item Compute for each scenario: $P_0$, Cohen's $\kappa$, weighted $\kappa$, Fleiss' $\kappa$, PABAK, positive/negative agreement, and confidence intervals.
\end{itemize}

\subsection*{Threshold Derivation (Three Strategies)}
\begin{enumerate}[nosep, leftmargin=1.5em]
    \item \textbf{Decision-utility mapping.} Map observed coefficients to acceptability under explicit false-positive/false-negative cost models. Produce threshold functions $T(\text{prevalence}, K, R, \text{cost ratio})$.
    \item \textbf{ROC-style optimization.} Treat ``acceptable vs.\ unacceptable reliability'' as a classification target. Choose cut points that minimize weighted misclassification.
    \item \textbf{Empirical clustering.} Cluster coefficient vectors within each context stratum. Check whether natural groupings correspond to meaningful differences in decision error.
\end{enumerate}

\subsection*{Validation}
\begin{itemize}[nosep, leftmargin=1.5em]
    \item Held-out simulation regimes and external empirical datasets (clinical coding, imaging, survey coding)
    \item Direct comparison against Landis-Koch (1977) and Altman (1991) baselines
    \item Quantify how often static ladders produce paradoxical mislabeling (Feinstein \& Cicchetti, 1990)
\end{itemize}

\section*{Deliverables}

\begin{itemize}[nosep, leftmargin=1.5em]
    \item \textbf{Context-aware threshold functions} mapping PABAK and $\kappa$ to interpretive categories, conditioned on prevalence, bias, and study design
    \item \textbf{Empirical mislabeling rates} for existing static scales across realistic conditions
    \item \textbf{Practitioner reporting checklist:} the minimal set of statistics required for defensible interpretation (cell counts, prevalence/bias indices, positive/negative agreement, CIs; cf.\ Cicchetti \& Feinstein, 1990; Sim \& Wright, 2005)
    \item \textbf{One-page practitioner guide} translating the framework into accessible decision support
\end{itemize}

\vspace{4pt}
\noindent\textit{This is not a new agreement metric. It is an empirically validated interpretation framework that connects agreement labels to decision consequences rather than historical convention.}

\section*{Key References}

\small
\begin{itemize}[nosep, leftmargin=1.5em]
\item Altman, D.\ G.\ (1991). \textit{Practical Statistics for Medical Research}. Chapman \& Hall.
\item Brennan, R.\ L.\ \& Prediger, D.\ J.\ (1981). Coefficient kappa: Some uses, misuses, and alternatives. \textit{Educ.\ Psychol.\ Meas.}, 41(3), 687--699.
\item Byrt, T., Bishop, J., \& Carlin, J.\ B.\ (1993). Bias, prevalence and kappa. \textit{J.\ Clin.\ Epidemiol.}, 46(5), 423--429.
\item Chen, G.\ et al.\ (2009). Measuring agreement of administrative data with chart data using prevalence unadjusted and adjusted kappa. \textit{BMC Med.\ Res.\ Methodol.}, 9:5.
\item Cicchetti, D.\ V.\ \& Feinstein, A.\ R.\ (1990). High agreement but low kappa: II.\ Resolving the paradoxes. \textit{J.\ Clin.\ Epidemiol.}, 43(6), 551--558.
\item Cohen, J.\ (1960). A coefficient of agreement for nominal scales. \textit{Educ.\ Psychol.\ Meas.}, 20(1), 37--46.
\item Feinstein, A.\ R.\ \& Cicchetti, D.\ V.\ (1990). High agreement but low kappa: I.\ The problems of two paradoxes. \textit{J.\ Clin.\ Epidemiol.}, 43(6), 543--549.
\item Hoehler, F.\ K.\ (2000). Bias and prevalence effects on kappa viewed in terms of sensitivity and specificity. \textit{J.\ Clin.\ Epidemiol.}, 53(5), 499--503.
\item Landis, J.\ R.\ \& Koch, G.\ G.\ (1977). The measurement of observer agreement for categorical data. \textit{Biometrics}, 33(1), 159--174.
\item Sim, J.\ \& Wright, C.\ C.\ (2005). The kappa statistic in reliability studies. \textit{Phys.\ Ther.}, 85(3), 257--268.
\item Vanacore, A.\ \& Pellegrino, M.\ S.\ (2022). Robustness of $\kappa$-type coefficients for clinical agreement. \textit{Stat.\ Med.}, 41(11), 1986--2004.
\end{itemize}

\end{document}
